\subsection{The prescribed involution survives a nonmetabelian step}

The following is a further test of the exact finite-action question
in \S52.2, folio 744, scan 193. It starts with the particular
metabelian solution of A.22, including its specified conjugators.
It does not replace that solution by a different finite action.

\begin{statement}{A.23}{Proved}{A compatible lift through the class-two open pro-5 kernel}
Let $F=\widehat F(x,y)$ and let
\[
 K=\ker\bigl(F\longrightarrow C_3,\quad x,y\longmapsto1\bigr).
\]
Let $N$ be the free class-two pro-$5$ quotient of $K$, and put
$Q=F/\ker(K\to N)$. Then the image of the A.22 involution
on $Q/Q''$ lifts to a continuous automorphism $\Phi$
of $Q$ of exact order two, whose outer class also has
exact order two. There are $f,h\in Q'$ mapping
to the specified A.22 conjugators such that
\[
 \Phi(x)=x,\qquad \Phi(y)=f^{-1}yf,\qquad
 \Phi(xy)=h^{-1}xyh,\qquad
 f\Phi(f)=h\Phi(h)=1.
\]
The open subgroup $N$ has a free rank-ten $\mathbb Z_5$
module of class-two BCH coordinates, and
$Q''\cong\mathbb Z_5$ is nontrivial and central.

In particular these formulas give compatible lifts on every
\[
 Q_n=F/\overline{K^{5^n}\gamma_3(K)},\qquad
 |Q_n|=3\cdot5^{10n},\qquad Q_n''\cong C_{5^n}.
\]
Thus the first genuinely nonmetabelian finite test, of order
$29\,296\,875$ for $n=1$, gives a lift rather than an
obstruction. No lift to $F$ is asserted.
\end{statement}
\begin{proof}
The Schreier basis of $K$ is
\[
 a=x^3,\qquad b=yx^{-1},\qquad
 c_0=xyx^{-2},\qquad d=x^2y.
\]
For completeness, the class-two pro-$5$ group on these four
generators is the BCH group of
\[
 L=V\oplus\bigwedge\nolimits^2 V,\qquad
 V=\mathbb Z_5 A\oplus\mathbb Z_5 B
       \oplus\mathbb Z_5 C\oplus\mathbb Z_5 D,
\]
with central $\bigwedge^2V$ and product
$u*v=u+v+\frac12[u,v]$. Here $A=\log a$,
$B=\log b$, $C=\log c_0$, $D=\log d$.
This description can be checked without a completion
assumption: in any class-two pro-$5$ group, collection of
powers of four generators followed by their six central
commutators gives the indicated continuous map from
$\mathbb Z_5^{10}$. The class-two collection identities
give its homomorphism law. Conversely the displayed BCH
group is generated topologically by $\exp A,\ldots,\exp D$,
and their six commutators are independent. These two maps
give the required universal property. Division by $2$
is allowed in $\mathbb Z_5$.

Conjugation by $x$ acts through the Lie automorphism $T$
given by
\[
 \begin{aligned}
 T(A)&=A,&T(B)&=C,\\
 T(C)&=D-A+\tfrac12[A,D],&
 T(D)&=A+B+\tfrac12[A,B].
 \end{aligned}
\]
The formulas follow directly from
$x b x^{-1}=c_0$,
$x c_0x^{-1}=da^{-1}$ and
$x d x^{-1}=ab$. They give $T^3=\operatorname{Ad}(a)$.
Since $3$ is a $5$-adic unit,
\[
 z=a^{-1/3}x,\qquad
 J=\operatorname{Ad}(a^{-1/3})T
\]
have order three. Thus
$Q=\exp(L)\rtimes\langle z\rangle$, with action $J$,
and
$x=\exp(A/3)z$, $y=\exp B\,x$.
The action $J$ fixes $A$.

Use the following fixed coordinate order:
\[
 A,B,C,D,\ [A,B],[A,C],[A,D],[B,C],[B,D],[C,D].
\]
Define $f=\exp u$ and $h=\exp v$ by
\[
 \begin{aligned}
 u={}&\tfrac23(A+B-D)
       -\tfrac4{27}[A,B]-\tfrac1{27}[A,C]
       +\tfrac5{27}[A,D]
       +\tfrac19[B,C]-\tfrac19[C,D],\\
 v={}&\tfrac23(A+C-D)
       +\tfrac1{27}[A,B]-\tfrac8{27}[A,C]
       +\tfrac1{27}[A,D]
       -\tfrac19[B,C]-\tfrac19[B,D].
 \end{aligned}
\]
All these coefficients lie in $\mathbb Z[1/3]\subset
\mathbb Z_5$. In the coinvariants of $J$ on $V$ one has
$B=C$ and $D=A+B$. The linear parts of $u,v$ therefore
vanish in the abelianization of $Q$, and every bracket
already lies in $Q'$. Consequently $f,h\in Q'$.

Substitution $x\mapsto x$, $y\mapsto f^{-1}yf$ gives
an endomorphism of $Q$. Indeed, lift $f$ to $K$.
The corresponding continuous endomorphism of $F$
preserves the map to $C_3$, hence $K$, and hence the
characteristic kernel of the free class-two pro-$5$
quotient of $K$. This proves descent; no invariance of
an arbitrary finite quotient is assumed.

Its action $P$ on $L$ fixes $A$ and sends the remaining
free Lie generators to
\[
 \begin{aligned}
 P(B)={}&-\tfrac23A-\tfrac13B+\tfrac23C+\tfrac23D
              -\tfrac29[A,C]-\tfrac19[A,D],\\
 P(C)={}&-\tfrac23A+\tfrac23B-\tfrac13C+\tfrac23D
              +\tfrac29[A,B]+\tfrac19[A,D],\\
 P(D)={}& \tfrac43A+\tfrac23B+\tfrac23C-\tfrac13D
              +\tfrac19[A,B]-\tfrac19[A,C].
 \end{aligned}
\]
On brackets $P([U,V])=[P(U),P(V)]$.
These formulas give, by substitution in the four free
generators and six brackets,
\[
 P^2=1,\qquad PJ=JP,\qquad P(A)=A,\qquad
 u*P(u)=v*P(v)=0.
\]
They also give the exact group identity
\[
              x f^{-1}y f=h^{-1}(xy)h.
\]
All identities hold over $\mathbb Q$ with coefficients
integral at $5$, hence hold in the full $5$-adic group.
They prove both norm equations and $\Phi^2=1$.
The linear part of $P(B)$ differs from $B$, so the order
is exactly two. The complete substitution checks are
also recorded in the executable certificate below.

We verify that this is a lift of the \emph{specified}
A.22 action. The second derived Lie module is the line
generated by
\[
 w=-[A,B]+[A,C]+[B,C]-[B,D]+[C,D].
\]
Indeed, the nontrivial $C_3$-character submodule of $V$
has basis $B-C,\ C-D+A$, whose bracket is $w$.
The full bracket module lies in $Q'$, so
$Q''=\exp(\mathbb Z_5w)$; $Jw=w$ and $w$ is central.
In particular this quotient is not metabelian.

On the abelian module $Q'/Q''$, write $X,Y,Z$ for
conjugation by $x,y,xy$, respectively, and put
\[
 E=1-\tfrac13(1+J+J^2).
\]
On this module, $J$ commutes with conjugation by every
$b\in N$: indeed $J(b)b^{-1}\in Q'$, which acts
trivially on $Q'/Q''$. Since
$x=\exp(A/3)z$ and $y=\exp B\,\exp(A/3)z$,
their commuting pro-$5$ factors contribute only
pro-$5$ actions, and both prime-to-$5$ parts are
$J$. Thus $E$ is exactly the image of the continuous
idempotent $e$ in A.22.
If $c=[x,y]$, the displayed $u,v$ satisfy, modulo
$\mathbb Z_5w$,
\[
 \begin{gathered}
 Eu=u,\qquad Ev=v,\qquad
 (X-1)(Y-1)u=-2E\log c,\\
 (Z-1)v=X(Y-1)u.
 \end{gathered}
\]
On the $E$-summand, $X-1,Y-1,Z-1$ are units, as in
A.22. These equations uniquely identify the images of
$f,h$ as its prescribed $c^r,c^s$. They establish the
claimed compatibility, rather than just the existence
of another involution on the new quotient.

The map to the order-$75$ certificate sends
$A\mapsto0$, $B\mapsto1$, $C\mapsto\zeta$,
$D\mapsto\zeta^2$ in $\mathbb F_{25,\mathrm{add}}$,
and kills all brackets. The formulas reduce to
$f=(3+4\zeta,0)$, $h=(4+3\zeta,0)$ and
$\bar\Phi(t,i)=(-t,i)$.
The generator identities prove that the quotient is
preserved. Its noninner action from A.22 proves that
$\Phi$ is outer.

Finally $5^nL$ is an invariant Lie ideal and BCH
subgroup. Reduction gives the stated $Q_n$, with
free rank-ten coordinate module over
$\mathbb Z/5^n$; $w$ remains primitive, so the second
derived subgroup has order $5^n$.
All reductions are compatible.

The program
\texttt{code/a\_metabelian\_lift\_class2.py}
performs two exact calculations. Its default run uses
$\mathbb F_5$ and checks the entire finite Lie basis
and all defining equations. Its
\texttt{--rational} run uses exact rational arithmetic,
checks the same identities, and verifies that every
coefficient is $5$-integral. The calculation solves
complete linear systems and checks all basis bracket
identities. It does not claim to enumerate all
$29\,296\,875$ group elements, and uses no sampling.
\end{proof}

\paragraph{The next relative quotient after A.23.}
The prescribed A.22 solution therefore has no obstruction
at the class-two open pro-$5$ stage. The next test is
the class-three quotient of the same $K$, followed by
class four, where the BCH denominators are still
$5$-adic units. A.24 below computes both complete
lifting systems. These relative pro-$5$ tests are
additional necessary stages; surviving them does not
prove a lift to the full free profinite group.
In particular the action on the two nontrivial cubic
character directions is $-1$ and its action on the
two invariant directions is $+1$. That homological
description supplies no geometric realization on
the three-punctured elliptic cover.

\begin{statement}{A.24}{Proved}{The full class-three and class-four lifting families}
Let $N_d$ be the free pro-$5$ quotient of $K$ of
nilpotency class at most $d$, and put
$Q_d=F/\ker(K\to N_d)$, for $d=3,4$.
Let $\mathcal X_d$ be the set of pairs of conjugators
$(f,h)$ which reduce to the A.23 pair in $Q_2$,
have the specified A.22 images in $Q_d/Q_d''$,
and satisfy the two norm and three peripheral
conditions in A.23.
Then $\mathcal X_3$ and $\mathcal X_4$ are nonempty
integral affine spaces:
\[
             \mathcal X_3\cong\mathbb Z_5^4,\qquad
             \mathcal X_4\cong\mathbb Z_5^{20}.
\]
Every pair gives a continuous automorphism of exact
order two whose outer class has exact order two.
These are counts of conjugator pairs, not counts of
distinct automorphisms.

The map $\mathcal X_4\to\mathcal X_3$ has image a
single point. In the standard Lyndon Lie basis with
$A<B<C<D$, four affine coordinates on $\mathcal X_3$
are the coefficients of
$[ADC],[ADD],[CCD],[CDD]$ in $\log h$.
The unique class-three point which extends is
\[
 \left(-\frac{11}{486},-\frac1{162},
              \frac{11}{486},-\frac{11}{486}\right).
                                                               \tag{A--III}
\]
In particular the initial class-three solution with
these four coordinates zero does not extend to
class four. Allowing the complete earlier family
removes that particular obstruction.

The coordinate ranks of $N_3,N_4$ are $30,90$,
and those of $Q_3'',Q_4''$ are $13,62$.
Reduction modulo $5^n$ gives $5^{4n}$ and
$5^{20n}$ conjugator pairs, respectively, above the
fixed A.23 pair; exactly one of the class-three
pairs extends to class four at that modulus.
The finite groups themselves have orders
$3\cdot5^{30n}$ and $3\cdot5^{90n}$.
The full free profinite lifting problem remains undecided.
\end{statement}
\begin{proof}
We state the complete equations and their exact finite
certificates, including the full solution spaces.
The calculation uses the free nilpotent Lie lattice,
not an arbitrary rational replacement of the pro-$5$
group. For class less than $5$, the exponential and
logarithm construction identifies the free nilpotent
$\mathbb Z_5$ Lie algebra with the free nilpotent
pro-$5$ group and preserves homomorphisms and normal
subgroups. A precise primary restatement is
\citet[\S2, Theorem 1(a)--(f), and the preceding free
object construction]{A_GaraialdeGonzalez2014}.
Its hypothesis is satisfied for $d=3,4$.

Let $L_d$ be the free Lie algebra on $A,B,C,D$,
truncated above degree $d$. Its homogeneous ranks are
$4,6,20,60$ through degree four. We compute inside
the free associative algebra truncated in the same
degrees, with
\[
 \exp u=\sum_{j=0}^{d}\frac{u^j}{j!},\qquad
 \log(1+u)=\sum_{j=1}^{d}\frac{(-1)^{j+1}u^j}{j}.
\]
The denominators are $5$-adic units. A Lyndon word
$w$ is bracketed recursively as
$[w]=[[a],[b]]$, where $b$ is its longest proper
Lyndon suffix and $w=ab$. Thus, for example,
$[ADC]=[[A,D],C]$ and $[CCD]=[C,[C,D]]$.
Ordering first by degree and then lexicographically
gives the integral basis used by the certificates.

The Schreier formulas, the automorphism $T$,
$z=a^{-1/3}x$ and the order-three action $J$ are
the exact group formulas in A.23, now evaluated
through degree $d$. In particular
\[
 T(A)=A,\quad T(B)=C,\quad
 T(C)=\log(\exp D\,\exp(-A)),\quad
 T(D)=\log(\exp A\,\exp B).
\]
Conjugating $T$ by $\exp(-A/3)$ gives $J$.
The program checks $T^3=\operatorname{Ad}(\exp A)$,
$J^3=1$, and all four Schreier identities.
It follows that this is the required quotient
$Q_d=\exp(L_d)\rtimes C_3$.

Write $u_2,v_2$ for the A.23 logarithms and let
\[
 u=u_2+\alpha,\qquad v=v_2+\beta,\qquad
 \alpha,\beta\in\bigoplus_{j=3}^{d}L_j.
\]
This includes every pair above the fixed class-two
pair. All such logarithms lie in $Q_d'$.
Let $\Phi_u$ be the endomorphism defined by
$x\mapsto x$, $y\mapsto\exp(-u)y\exp u$,
and let $P_u$ be its action on $L_d$, obtained
from the four Schreier generators.
The same characteristic-kernel argument as in
A.23 proves that the substitution descends.

Put $H=\log(Q_d')$ and $H_2=\log(Q_d'')$. Explicitly,
\[
 H=\langle B-C,\ C-D+A\rangle_{\mathbb Z_5}
                   +\bigoplus_{j=2}^{d}L_j,\qquad
 H_2=[H,H].
                                                               \tag{A--IV}
\]
The bracket span is an ideal: $H$ is an ideal and
Jacobi shows $[L_d,[H,H]]\subset[H,H]$.
Lazard's correspondence therefore identifies it
with the second derived subgroup. Its ranks are
computed as $13$ and $62$, respectively; the same
ranks modulo $5$ certify saturation of these
integral spans.

Let $X,Y,Z$ be conjugation by $x,y,xy$ on $L_d$,
and put $E=1-(1+J+J^2)/3$. As in A.23, the
induced operators on $H/H_2$ have prime-to-$5$
parts $J,J,J^2$. The following four linear
conditions are therefore exactly the required
A.22 metabelian residues:
\[
 \begin{gathered}
 Eu-u,\quad Ev-v,\quad
 (X-1)(Y-1)u+2E\log[x,y],\\
 (Z-1)v-X(Y-1)u
                       \quad\in H_2.
 \end{gathered}                                                  \tag{A--V}
\]
The units $X-1,Y-1,Z-1$ on the $E$-summand
make this identification unique.
The remaining three equations are exact group
equations, calculated with the displayed
exponential and logarithm:
\[
 \begin{aligned}
 &\log\bigl(\exp u\,\exp(P_u u)\bigr)=0,\\
 &\log\bigl(\exp v\,\exp(P_u v)\bigr)=0,\\
 &x\exp(-u)y\exp u
             =\exp(-v)(xy)\exp v.
 \end{aligned}                                                   \tag{A--VI}
\]
Together (A--V)--(A--VI) are necessary and sufficient
for the stated pair problem.

All seven residuals are affine in the coefficients
of $\alpha,\beta$ through degree four. To see this,
each correction has degree at least three.
In a direct product two corrections first interact
in degree at least six. In substitution into a
degree-three correction, a changed generator has
degree at least three and therefore first produces
a quadratic term in degree at least five.
There are no such terms in $L_3$ or $L_4$.
The complete linear systems are consequently
obtained by evaluating the residual at zero and
at every single basis correction.

Here are the exact ranks. The row counts retain
zero and dependent equations, rather than deleting
them by an unspecified preprocessing step.
\[
\begin{array}{c|r|r|r|r|r}
d&\operatorname{rank}L_d&\operatorname{rank}H_2&
 \text{system size}&\text{rank}&\text{nullity}\\ \hline
3&30&13&140\times40&36&4\\
4&90&62&560\times160&140&20
\end{array}
\]
The program
\texttt{code/a\_metabelian\_lift\_higher.py}
saves each sparse matrix, right-hand side, complete
rational particular solution, complete nullspace
basis, pivot columns, and free-variable coordinates.
Every particular-solution and nullspace coefficient
is $5$-integral. The free columns in each nullspace
basis form an identity matrix.

The independent verifier
\texttt{code/a\_higher\_lift\_verify.py} does not import
the BCH construction. It checks $Aq=b$ and $Av_i=0$
for every displayed basis vector, verifies the
identity free-column minor, and computes the ranks
$36,140$ independently modulo $5$.
The independent nullspace vectors give the rational
upper rank bounds, and these nonzero modular minors
give the matching lower bounds. It also checks
that reduction of the rational system is exactly
the separately constructed $\mathbb F_5$ system.
Thus the certificates prove the complete integral
affine solution spaces, with no undetermined
integrality or rank assumption.

For clarity, the saved full coefficient vectors
are in
\texttt{code/a\_higher\_lift\_coordinates.csv};
the full affine bases and equations remain in the
class-three and class-four JSON certificates.
Their exact hashes are recorded beside the witnesses.
This is an exact computer-assisted verification of
finite equations; no group elements or possible
solutions were sampled.

The restriction calculation is also performed on
the entire solution spaces. Every one of the
twenty class-four nullspace vectors has zero
degree-three part. Restricting its particular
solution to degree three gives a solution of the
class-three system with the four free coordinates
(A--III). Conversely that point has all the
displayed class-four lifts. This proves both
uniqueness of the extendible point and the failure
of the earlier zero-parameter choice. In
particular the class-four computation has not
silently frozen one arbitrary class-three lift.

Since the matrix ranks do not drop modulo $5$,
their nonzero Smith factors over $\mathbb Z_5$
are units. Reduction modulo $5^n$ therefore gives
precisely the asserted numbers of pairs, and the
constant restriction map remains constant.
The norms give $\Phi_u^2=1$; its nontrivial outer
image on the order-$75$ quotient proves that both
the actual and outer orders are exactly two.
\end{proof}

\paragraph{The first remaining integral stage.}
The next specified relative quotient has class five
in the open pro-$5$ kernel. The class-four family
has twenty free parameters, all in degree four.
They must be retained in that next lifting problem.
In degree five the rational BCH denominators can
contain $5$, so the bare lattice
$L_{\leq5}(\mathbb Z_5)$ with the same formulas
cannot simply be substituted for the free
class-five pro-$5$ group. A rational solution
with a denominator $5$ would not by itself prove
a pro-$5$ obstruction. The next calculation
requires integral Hall group coordinates, or an
equivalent explicit description of the correct
lattice, together with the complete preceding
parameter family. The separate full profinite
lifting obligation in A.12 also remains.
